\section{Limitations and Future Work}\label{sec:limitations}

\textbf{Hardware coverage is one kernel, one test, one measurement.} The Shor kernel's \texttt{m\_axi} fix is now hardware-verified (Section~\ref{subsec:res-hardware}): it works, and its post-route reports are real. That is one kernel out of three, one test (the 27-case self-test, not the full \(3^9\) exhaustive sweep), and one opportunistic, non-rigorous latency measurement. The repetition-code and Steane kernels have not been placed and routed or tested at all, and remain HLS-ESTIMATE throughout Section~\ref{sec:results}. Extending hardware coverage to them, running the exhaustive \(3^9\)/\(3^7\) sweeps, and replacing the opportunistic latency number with a rigorous C++-hosted, tail-resolving measurement (ideally pinned to an isolated core under \texttt{SCHED\_FIFO}, at \(\ge10^6\) shots) are the immediate next steps, blocked only on continued access to a card shared with other work on this machine.

\textbf{Latency measurement is preliminary.} The 49.5\,\textmu{}s median / 61.0\,\textmu{}s p99 round trip reported in Section~\ref{subsec:res-hardware} used a single Python process, no core isolation, no real-time scheduling, 10,000 shots (not the \(\ge10^6\) a tail distribution needs), and did not decompose the round trip into its PCIe, AXI-Lite, and host-software components. It establishes that host dispatch, not kernel logic, dominates observed latency by roughly three orders of magnitude, which is directionally consistent with what a real-time feasibility assessment would need to know, but it is not that assessment.

\textbf{Steane kernel is a reconstruction.} The three-mode Steane kernel evaluated in this paper was rebuilt from a textual description because the original source could not be located (Appendix~\ref{app:reconstruction}). Every Steane number in this paper should be read as a characterisation of this reconstruction, not a verification of whatever kernel actually produced a prior submission's Steane figures.

\textbf{Single-code circuit-level coverage.} The only circuit-level, multi-round result in this paper (Section~\ref{subsec:res-circuit}) covers the repetition code. Building equivalent Stim circuits for the Shor and Steane codes, and swapping in this project's own decoder logic rather than PyMatching's general MWPM, is real engineering work not attempted here.

\textbf{No batching, no multi-compute-unit measurement.} This revision does not report a batched-throughput sweep, a multi-compute-unit resource or throughput study, or any fair CPU/GPU baseline comparison. A batched, HBM-backed kernel architecture exists in the accompanying repository (unrelated to the kernels evaluated in this paper) as a candidate starting point for that work.

\textbf{Distance-5 and beyond.} As Table~\ref{tab:lut-scaling} makes explicit, LUT-based decoding does not extend past roughly 16 syndrome bits. A distance-5 rotated surface code with a union-find decoder is a plausible next step and is scoped, with a time estimate, in the accompanying decision record, but has not been started.

\textbf{Noise models.} Every result in this paper uses either a phenomenological Pauli-injection model or, for the repetition code only, a uniform circuit-level depolarising model. Neither captures leakage, correlated errors, or noise calibrated against a real device's measured error rates, all of which a deployment-relevant assessment would need.
